\newpage
\appendix
\section*{Supplementary Information}

\setcounter{section}{0}
\renewcommand{\thesection}{S\arabic{section}}
\renewcommand{\theequation}{S\arabic{equation}}
\setcounter{equation}{0}
\renewcommand{\thefigure}{S\arabic{figure}}
\setcounter{figure}{0}
\renewcommand{\thetable}{S\arabic{table}}
\setcounter{table}{0}

\section{Cosserat Rod Geometry and Kinematic Parameterization}

We model the human spine as a Cosserat rod, defined by a centerline curve $\mathbf{r}(s,t) \in \mathbb{R}^3$ parameterized by arc length $s \in [0, L]$. This dimensionality reduction from 3D elasticity to a 1D director theory is justified by the high aspect ratio of the vertebral column ($L/d \approx 13$). To each point $s$, we attach an orthonormal material frame $\{\mathbf{d}_1(s,t), \mathbf{d}_2(s,t), \mathbf{d}_3(s,t)\}$ describing the cross-sectional orientation~\cite{goriely2017mathematics, gazzola2018forward}. The kinematics are governed by:
\begin{align}
\mathbf{r}'(s) &= \mathbf{v}(s) = \mathbf{d}_3(s) + \boldsymbol{\varepsilon}(s), \\
\mathbf{d}_i'(s) &= \mathbf{u}(s) \times \mathbf{d}_i(s),
\end{align}
where $\mathbf{v}(s)$ is the tangent vector, $\boldsymbol{\varepsilon}(s)$ represents shear and extension strains, and $\mathbf{u}(s) = \kappa_1 \mathbf{d}_1 + \kappa_2 \mathbf{d}_2 + \tau \mathbf{d}_3$ is the Darboux vector encoding flexural curvature $\kappa_{1,2}$ and torsional twist $\tau$.

\paragraph{Boundary Conditions.} The sacral base ($s=0$) is clamped and the cranial tip ($s=L$) is free:
\begin{align}
\mathbf{r}(0) = \mathbf{0}, \quad &\mathbf{d}_i(0) = \mathbf{e}_i \quad (\text{Clamped Base}) \nonumber \\
\mathbf{n}(L) = \mathbf{0}, \quad &\mathbf{m}(L) = \mathbf{0} \quad (\text{Free Tip})
\end{align}

\section{Cosserat Force and Moment Balance}

For a static rod subject to gravity $\mathbf{f}_g = \rho A \mathbf{g}$ and active moments induced by the information field, the balance equations for internal force $\mathbf{n}$ and moment $\mathbf{m}$ are:
\begin{align}
\mathbf{n}'(s) + \mathbf{f}_g &= \mathbf{0}, \nonumber \\
\mathbf{m}'(s) + \mathbf{r}'(s) \times \mathbf{n}(s) + \mathbf{m}_{\mathrm{info}}'(s) &= \mathbf{0},
\end{align}
where $\mathbf{m}_{\mathrm{info}}$ represents the active couple generated by the information field gradient.

\section{Developmental Information Field Parameterization}

The scalar information field $I(s)$ is parameterized as a bimodal Gaussian:
\begin{equation}
I(s) = A_c \exp\left[-\frac{(s/L - s_c)^2}{2\sigma_c^2}\right] + A_l \exp\left[-\frac{(s/L - s_l)^2}{2\sigma_l^2}\right] + I_0,
\end{equation}
where $A_c = 0.5$, $s_c = 0.80$, $\sigma_c = 0.08$ (cervical lordosis), $A_l = 0.7$, $s_l = 0.25$, $\sigma_l = 0.10$ (lumbar lordosis), and $I_0 = 0.3$ (baseline).

The information content is quantified using the local Shannon entropy:
\begin{equation}
\ShannonH[I] = -\int I(s) \log I(s) \, ds,
\end{equation}
which bounds the precision of the target metric specification.

\section{Biological Metric and Information Geometry}

The central hypothesis of the IEC framework is that developmental information modifies the effective metric experienced by the rod's reference manifold. We define a biological metric:
\begin{equation}
d\ell_{\mathrm{eff}}^2 = g_{\mathrm{eff}}(s)\,ds^2 = \exp\left[2\left(\beta_1 \tilde{I}(s) + \beta_2 \frac{\partial \tilde{I}}{\partial s}\right)\right] ds^2,
\end{equation}
where $\tilde{I}$ is the normalized information field and $\beta_{1,2}$ are dimensionless coupling constants.

\section{Geodesic Deviation and Regime Classification}

We quantify the degree of biological counter-curvature using the Kullback--Leibler divergence between the IEC-coupled geometry and the passive gravitational state:
\begin{equation}
\KLDivergence = \int_0^L \left[ \frac{1}{2} \left( \kappa_{\mathrm{IEC}}(s) - \kappa_{\mathrm{passive}}(s) \right)^2 \FisherInfo(s) \right] w(s)\, ds,
\end{equation}
where $\FisherInfo(s)$ is the local Fisher Information and $w(s) = g_{\mathrm{eff}}(s)$. The normalized geodesic deviation is:
\begin{equation}
\Dgeohat = \frac{\KLDivergence}{\KLDivergence^{\mathrm{max}}(\chi_\kappa, g)}.
\end{equation}

Regime thresholds: gravity-dominated ($\Dgeohat < 0.1$), cooperative ($0.1 \leq \Dgeohat \leq 0.3$), information-dominated ($\Dgeohat > 0.3$). These were determined by analyzing the bifurcation of the lowest eigenvalue $\lambda_0$.

\section{Metabolic Power Density and Dissipation}

The metabolic power density required to sustain counter-curvature is:
\begin{equation}
\dot{\mathcal{F}} = \int_0^L \left[ \eta_p \left|\frac{\partial \kappa}{\partial t}\right|^2 + \eta_a \left(\kappa(s) - \kappa_{\mathrm{passive}}(s)\right)^2 + \Gamma_m(s) \right] ds,
\end{equation}
where $\eta_p$ is the proprioceptive feedback dissipation coefficient, $\eta_a$ is the active moment maintenance coefficient, and $\Gamma_m(s)$ is the basal tissue maintenance cost density.

\section{Protein Dynamics and Energy Competition}

During adolescent growth, mechanosensory and growth/maintenance protein classes compete for a finite metabolic energy pool:
\begin{equation}
E_{total}(t) = E_{mechano}(t) + E_{growth}(t) + E_{metabolic}
\end{equation}
\begin{equation}
\frac{dE_{mechano}}{dt} = k_{syn,m} \cdot [Resources] - k_{deg,m} \cdot [Mechano] - \dot{\mathcal{F}}_{load}
\end{equation}
where $\dot{\mathcal{F}}_{load} \propto L^4$ represents the increasing thermodynamic load. The growth protein synthesis rate is driven by GH/IGF-1 and scales with $L^2$ (surface area of nutrient diffusion), creating the scaling mismatch at the molecular level.

The vulnerability is quantified via:
\begin{equation}
\Lambda(t) = \frac{dV/dt}{\tau_{adapt}}
\end{equation}
where $dV/dt$ is volumetric growth rate and $\tau_{adapt}$ is the mechanosensory adaptation time constant.

\section{The Spinal Clock and Convective Shutdown}

The intervertebral disc possesses an autonomous circadian clock driven by the BMAL1/CLOCK transcriptional loop~\cite{dudek2017clock}. Gravity provides the synchronization signal through daily loading cycles: daytime compression forces fluid out of the nucleus pulposus, while nocturnal unloading draws fluid back via osmotic pressure~\cite{urban2004nutrition}.

Under microgravity or prolonged bed rest, this pump arrests (``Convective Shutdown''), creating hypoxia that stabilizes HIF-1$\alpha$ and upregulates catabolic enzymes (MMP1, MMP3). The resulting matrix degradation reduces $B_{\mathrm{eff}}(s)$ and widens the Energy Deficit Window.

This mechanism generates three falsifiable predictions:
\begin{enumerate}
    \item Clock gene desynchronization (Per2, Bmal1 flattening in IVD) within 72 hours of unloading, preceding structural changes.
    \item HIF-1$\alpha$ stabilization in the nucleus pulposus preceding scoliotic curvature; HIF-1$\alpha$ inhibitors should delay onset.
    \item Pulsatile axial compression at the appropriate circadian phase will preserve spinal geometry more effectively than static compression.
\end{enumerate}

\section{Complete Theoretical Derivations}
\label{sec:supp_theory}

\textit{This section provides the complete mathematical derivations, delay differential equation analysis, Hopf bifurcation calculations, and detailed coupling mechanism formulations supporting the main-text conceptual summary in Section 2.}

\subsection{The Bio-Gravitational Number and the Allometric Trap}

For a vertebral column of length $L$, flexural stiffness $EI$, mass $M$, and gravitational acceleration $g$, we define the \textbf{Bio-Gravitational Number}:
\begin{equation}
\mathcal{B}_g = \frac{EI}{M g L^2},
\label{eq:bg_number_passive_supp}
\end{equation}
a dimensionless ratio of passive flexural resistance to gravitational bending moment. When $\mathcal{B}_g > 0.1$, an organism can resist gravitational collapse through passive stiffness alone. When $\mathcal{B}_g < 0.1$, the organism enters the ``Allometric Trap'': it must actively maintain its shape at continuous metabolic cost. Cross-species analysis reveals that small quadrupeds operate well above this threshold, while humans ($\mathcal{B}_g \approx 0.01$) are deep within the trap, relying almost entirely on active metabolic maintenance to prevent spinal collapse~\cite{smit2020adolescent}. The Bio-Gravitational Number is analogous to the Froude number for locomotion, defining a fundamental physical boundary for spinal stability. However, scaling exponents vary systematically depending on physiological state~\cite{glazier2022variable}, and vertebral scaling exhibits size-dependent breakpoints~\cite{smith2024small}, emphasizing that bipedalism and human-specific axial loading uniquely exacerbate this vulnerability~\cite{gorman2009idiopathic}.

The central problem of vertebrate morphogenesis is the translation of a one-dimensional genetic code into a robust three-dimensional geometry that can withstand environmental loads. Classical structural mechanics predicts that a flexible column clamped at the base and subjected to its own weight will buckle or sag into a monotonic C-shape~\cite{goriely2017mathematics}. Yet, the healthy human spine exhibits a complex, alternating S-shape (lordosis and kyphosis).

This anomaly suggests that the S-shape is not a passive equilibrium but an active \emph{Counter-Curvature}---a geometric standing wave maintained by the continuous expenditure of metabolic energy against the gravitational field. This maintenance constitutes a biological implementation of \emph{optimal feedback control}~\cite{todorov2002optimal} and can be conceptually modeled using \emph{Latent Imagination}~\cite{hafner2019dreamer}, where the nervous system learns an internal world model of gravity and biomechanics to predict and optimize long-horizon postural behaviors by minimizing a cost function comprising both ``geometric error'' (deviation from the target shape) and ``metabolic effort'' (ATP consumption).

\subsection{Delay Differential Equations (DDE) and the Derivative Gain Trap}

The active maintenance of spinal alignment relies on proprioceptive feedback from paraspinal muscles and mechanosensors. However, neural conduction and central processing are not instantaneous. Established models of human postural control demonstrate that the upright state is stabilized by proportional-derivative (PD) or proportional-derivative-acceleration (PDA) feedback, subject to a time delay $\tau$~\cite{milton2009time, insperger2013acceleration}. We formalize the spine as an active control system governed by a Delay Differential Equation (DDE).

The proprioceptive control law generating corrective force $F_{control}$ with feedback delay $\tau$ is:
\begin{equation}
\boldsymbol{\kappa}_{\mathrm{rest}}(s) = \boldsymbol{\kappa}_{\mathrm{gen}} + \chi_\kappa \nabla I(s).
\label{eq:kappa_rest_field_supp}
\end{equation}
where $K_P$ and $K_D$ are proportional and derivative gains, and the error is $e(s,t) = \kappa_{measured}(s,t) - \kappa_{target}$.

Linearizing around the straight configuration, the dynamics of the spinal column become:
\begin{equation}
\rho A \frac{\partial^2 \kappa}{\partial t^2} + \eta \frac{\partial \kappa}{\partial t} + EI \frac{\partial^4 \kappa}{\partial s^4} + K_P \kappa(s, t-\tau) + K_D \frac{\partial \kappa}{\partial t}(s, t-\tau) = 0
\label{eq:dde_spine_supp}
\end{equation}

The stability of this active structure is governed by the dimensionless Bio-Gravitational Number $\mathcal{B}_g$:
\begin{equation}
\mathcal{B}_g = \frac{\chi_M \langle |\nabla I| \rangle}{\rho A g L^2}.
\label{eq:bg_number_info_supp}
\end{equation}
Crucially, when $\mathcal{B}_g > 1$, the biological system effectively dominates gravitational mechanics. At the molecular level, this coupling relies on the structural rigidity of patterning transcription factors; AlphaFold structural analysis of critical HOX proteins (e.g., HOXC8, UniProt P31273; HOXB13, UniProt Q92826) reveals conserved, high-confidence DNA-binding domains capable of enforcing the sharp spatial identity boundaries required by our field parameterization.

During the adolescent growth spurt, the spine elongates rapidly. Because nerve conduction velocity in the extremities does not scale proportionately with rapid height increases~\cite{lang1985nerve, robinson2001nerve}, the absolute neural feedback delay $\tau$ increases significantly. We term this vulnerability the \textbf{Derivative Gain Trap}: derivative gains ($K_D$) that provide essential damping and stability during childhood become destabilizing as the delay $\tau$ crosses a critical threshold $\tau_c$.

\subsection{Hopf Bifurcation and the Onset of Scoliosis}

When the neural delay exceeds the critical threshold ($\tau > \tau_c$), the system undergoes a supercritical Hopf bifurcation~\cite{landry2005dynamics}. The previously stable upright equilibrium becomes unstable, and stable limit cycles (oscillations) emerge. In the context of the spine, these delay-induced oscillations manifest as continuous, micro-mechanical buckling events that asymmetric biomechanical loading then locks into a permanent scoliotic deformity.

Assuming solutions of the form $\kappa(s,t) = A e^{\lambda t} \sin\left(\frac{n\pi s}{L}\right)$, the characteristic equation is:
\begin{equation}
\boldsymbol{\kappa}_{\mathrm{rest}}(s) = \boldsymbol{\kappa}_{\mathrm{gen}} + \chi_\kappa \nabla I(s),
\label{eq:kappa_rest_iec_supp}
\end{equation}
where $\chi_\kappa$ is a geometric coupling constant and $\boldsymbol{\kappa}_{\mathrm{gen}}$ is the baseline genetic curvature.

\paragraph{Stiffness modulation (IEC-2).}
The information field modulates the effective bending and torsional stiffness:
\begin{equation}
B_{\mathrm{eff}}(s) = E_0 I_{\mathrm{area}} (1 + \chi_E I(s)), \qquad
C_{\mathrm{eff}}(s) = G_0 J (1 + \chi_C I(s)).
\label{eq:stiffness_modulation_supp}
\end{equation}

The total energy functional governing the rod's equilibrium is:
\begin{equation}
\mathcal{E}_{\mathrm{total}} = \int_0^L \left[ \frac{1}{2} B_{\mathrm{eff}}(\kappa - \kappa_{\mathrm{rest}})^2 + \frac{1}{2} C_{\mathrm{eff}} \tau^2 + \frac{1}{2} K_{\mathrm{eff}} \varepsilon^2 + \rho A \mathbf{r} \cdot \mathbf{g} \right] ds.
\label{eq:iec_energy_supp}
\end{equation}

For purely imaginary roots $\lambda = \pm i\omega$ (defining the bifurcation boundary), the system transitions from damped stability to active oscillation. This mathematical control theory precisely explains why Adolescent Idiopathic Scoliosis (AIS) uniquely coincides with the rapid adolescent growth spurt: it is the exact developmental window where increasing $\tau$ pushes the neuromotor control system across the Hopf bifurcation boundary.

\subsection{The Energy Deficit Window: NAD+ and Proprioceptive Integrity}

The biological spine is a dissipative structure requiring continuous metabolic expenditure to maintain its sensory and structural integrity. The metabolic power required to sustain this configuration scales superlinearly with length ($L^4$ demand vs $L^2$ surface-limited supply)~\cite{sato2025convective}.

This scaling mismatch creates a transient \textbf{Energy Deficit Window} during the adolescent growth spurt. We propose that this metabolic deficit directly modulates the control parameters in the DDE model. Specifically, proprioceptive Ia/II afferents are among the largest and most metabolically demanding neurons.

Recent evidence demonstrates that NAD+ depletion during development causes severe spinal deformities~\cite{basse2021nampt}, while oxidative stress directly drives scoliosis in animal models~\cite{pumputis2025oxidative}. Furthermore, NAD+ insufficiency activates the SARM1 pathway, which triggers local axon degeneration~\cite{gerdts2015sarm1}. We hypothesize that the adolescent Energy Deficit Window causes localized NAD+ insufficiency, which mildly compromises proprioceptive axon integrity. This subclinical neuropathy further increases the effective neural delay $\tau$, accelerating the transition across the Hopf bifurcation boundary and triggering scoliotic buckling.

\subsection{Molecular Grounding: The AlphaFold Anisotropy Gap (Exploratory)}

While the primary failure mode is governed by macroscopic delay dynamics, we also explore the molecular architecture of the mechanosensory apparatus. AlphaFold v6 structural analysis of 23 key spinal proteins reveals that demand-side mechanosensory proteins (e.g., Vimentin, PIEZO2) exhibit significantly higher structural anisotropy (mean $3.08 \pm 1.44$) than supply-side metabolic regulators (mean $1.79 \pm 0.56$).

We tentatively propose a \textbf{Demand-Supply Anisotropy Gap} ($p = 0.011$ in our dataset): the proteins required to \emph{sense} mechanical geometry are structurally extended and thermodynamically expensive to maintain, while their metabolic regulators are compact and often intrinsically disordered. We emphasize that this observation is hypothesis-generating. Current AI structural prediction tools cannot directly infer mechanical properties under load~\cite{gomes2022protein, zonta2024sequence}, and confidence scores for extracellular matrix components are frequently low. Nonetheless, this structural asymmetry provides a compelling direction for future experimental validation of the cellular mechanisms underlying the Energy Deficit Window.


\section{Detailed Methods}
\label{sec:supp_methods}

\textit{This section provides detailed computational protocols, steered molecular dynamics parameters, tissue homogenization calculations, and protein dynamics modeling supporting the main-text Methods summary in Section 3.}

\subsection{Multi-Scale Protein Mechanics Pipeline}

To ground the macroscopic parameters $\eta_p, \eta_a, \Gamma_m$ in molecular reality, we developed a multi-scale pipeline integrating AlphaFold predictions with structural energetics.
\begin{enumerate}
    \item \textbf{Structure Prediction:} We retrieved high-confidence structures for 23 key spinal proteins from the AlphaFold Protein Structure Database.
    \item \textbf{Structural Metrics:} For each protein, we computed the \emph{Anisotropy Index} (ratio of principal moments of inertia) and \emph{Radius of Gyration} ($R_g$) using the `afcc` module. These metrics serve as proxies for the entropic cost of maintaining the protein's orientation against thermal noise.
    \item \textbf{Energetic Mapping:} We defined the metabolic cost of maintenance $C_{maint} \propto \text{Anisotropy} \times N_{residues}$. This assumes that extended, anisotropic structures require more frequent chaperone intervention and cytoskeletal tension to prevent misfolding or aggregation compared to globular proteins.
    \item \textbf{Steered Molecular Dynamics (SMD):} SMD simulations were executed in GROMACS using CHARMM36m and TIP3P water. We applied constant-velocity pulling to extract force-extension curves, enabling calculation of single-molecule stiffness $k_{molecule}$ from the low-strain linear regime.
    \item \textbf{Fibril Assembly and Mechanics:} Fibrillar assemblies were coarse-grained with explicit D-band organization and cross-link density. Effective fibril moduli were computed incorporating collagen axial mechanics and proteoglycan osmotic compression terms.
    \item \textbf{Tissue Homogenization:} Regional tissue elasticity tensors were computed using mixture theory: $\mathbf{C}_{tissue}(s) = \sum \phi_i(s) \mathbf{R}(\theta_i) \mathbf{C}_i \mathbf{R}(\theta_i)^T$, weighting constituent stiffness $\mathbf{C}_i$ by volume fractions $\phi_i$ from histology and orientation $\theta_i$ from polarized imaging.
    \item \textbf{Validation:} We verified predicted molecular stiffness against AFM nanoindentation measurements and bulk mechanical tissue tests, achieving an $R^2 \approx 0.81$.
\end{enumerate}

\subsection{Protein Dynamics and Energy Modeling}

We simulated the temporal competition between protein classes using a system of coupled differential equations. The model tracks the energy pool $E(t)$ available for protein synthesis/maintenance over the 5--20 year age range.
\begin{itemize}
    \item \textbf{Supply:} Modeled as $S(t) = S_0 (L(t)/L_0)^2$, reflecting the diffusion-limited transport through the vertebral endplate surface area.
    \item \textbf{Demand:} Modeled as $D(t) = D_{maint} (L/L_0) + D_{grav} (L/L_0)^4$, capturing the linear scaling of tissue volume maintenance and the quartic scaling of gravitational moment opposition.
    \item \textbf{Dynamics:} The fidelity of the mechanosensory array $F(t)$ evolves according to $dF/dt = k_{repair}(S - D)$ when $S > D$, and $dF/dt = k_{damage}(S - D)$ when $S < D$ (deficit), bounded to $[0, 1]$.
\end{itemize}

